\chapter{Vaginal bleeding}
\index{Vaginal bleeding}

\section*{Definition and Overview}
Vaginal bleeding is bleeding from the vagina or the organs behind it, including the uterus, cervix, and surrounding tissues. It can range from light spotting to heavy, life-threatening loss, and it occurs in several distinct contexts: normal menstruation, abnormal bleeding between or after periods, bleeding after sexual intercourse, bleeding after menopause, and bleeding in pregnancy. The term abnormal vaginal bleeding covers bleeding that is heavier, longer, more frequent, or more irregular than a woman's usual pattern. In pregnancy, bleeding raises entirely different concerns, including miscarriage, ectopic pregnancy, and complications such as placental problems. Because the causes range from trivial to catastrophic, vaginal bleeding is always taken seriously, and the first priorities are to determine how heavy the bleeding is, whether the woman is pregnant, and whether she is haemodynamically stable. It is also a symptom that can cause great anxiety, so clear communication about what is happening and what to expect is part of good care.

\section*{Epidemiology}
Abnormal uterine bleeding is one of the most common reasons for gynaecological consultation, affecting around 1 in 3 women at some point in their lives and accounting for up to a third of gynaecological outpatient visits. Heavy menstrual bleeding alone affects roughly 1 in 4 women of reproductive age. Bleeding in early pregnancy is very common, occurring in about 20--25\% of pregnancies, and about half of these continue normally. Around 10--15\% of clinically recognised pregnancies end in miscarriage. Bleeding after menopause occurs in approximately 5--10\% of postmenopausal women and, while most cases are benign, it is the classic symptom of endometrial cancer, which is why it always warrants investigation. Bleeding after intercourse affects about 1 in 20 women per year and is usually benign, but it can be a sign of cervical infection or, rarely, cancer. Bleeding associated with an intrauterine device or with the combined oral contraceptive is common, especially in the first months of use. The pattern of bleeding across a woman's life also matters: new bleeding at any age deserves review, but the urgency of the assessment is greatest in pregnancy and after menopause.

\section*{Aetiology and Pathophysiology}
Vaginal bleeding arises when the lining of the uterus or the surface of the cervix or vagina bleeds. The causes depend heavily on age, pregnancy status, and the pattern of bleeding.

\begin{itemize}
    \item \textbf{Physiological causes:} normal menstruation, implantation bleeding in early pregnancy, and bleeding around ovulation
    \item \textbf{Hormonal disturbances:} polycystic ovary syndrome, thyroid disease, perimenopause, stress, and weight change cause irregular, heavy, or prolonged bleeding
    \item \textbf{Structural causes:} fibroids, endometrial polyps, adenomyosis, and uterine or cervical abnormalities
    \item \textbf{Pregnancy-related causes:} threatened or actual miscarriage, ectopic pregnancy, and, later in pregnancy, placenta praevia or placental abruption
    \item \textbf{Infection:} cervicitis, pelvic inflammatory disease, sexually transmitted infections, and vaginitis can cause bleeding, especially after sex
    \item \textbf{Medications:} anticoagulants (blood thinners), some hormonal treatments, and tamoxifen
    \item \textbf{Contraception:} breakthrough bleeding with hormonal contraception, and bleeding associated with an intrauterine device (IUD)
    \item \textbf{Cancer:} cervical cancer and endometrial cancer, which is why postmenopausal bleeding and bleeding between periods need urgent review
    \item \textbf{Trauma:} injury, foreign bodies, sexual assault, or recent procedures
    \item \textbf{Coagulation and blood disorders:} inherited bleeding disorders, liver disease, and low platelet counts can cause heavy or prolonged bleeding
\end{itemize}

\section*{Clinical Features}
The features that matter most are the amount of bleeding, its timing, and any associated symptoms. Recording these details helps the clinician decide how urgent the problem is and what tests are needed.

\begin{itemize}
    \item Bleeding that soaks pads or tampons hourly, passes large clots, or leaks onto clothing
    \item Bleeding that lasts more than 7 days or occurs more often than every 21 days
    \item Spotting between periods or bleeding after sexual intercourse
    \item Any bleeding after menopause
    \item Bleeding in pregnancy, with or without pain
    \item Associated cramps, back pain, or a dull ache in the pelvis
    \item Dizziness, light-headedness, paleness, or fainting with heavy bleeding
    \item Fever, discharge, or an offensive smell, which may indicate infection
\end{itemize}

\section*{Differential Diagnosis}
\begin{itemize}
    \item Normal menstruation and variations in cycle length
    \item Miscarriage, ectopic pregnancy, and other pregnancy complications
    \item Uterine fibroids, polyps, and adenomyosis
    \item Hormonal disorders including PCOS and thyroid disease
    \item Sexually transmitted infections and pelvic inflammatory disease
    \item Cervical pathology, including cervical polyps, cervicitis, and cervical cancer
    \item Endometrial hyperplasia and endometrial cancer
    \item Bleeding disorders and anticoagulant medication
    \item Perimenopause and menopause
\end{itemize}

\section*{When to Seek Medical Care}
A doctor should be consulted for any unexpected bleeding, including bleeding between periods, bleeding after sex, heavier or more painful periods than usual, or any bleeding after menopause. All women who could be pregnant should seek advice for any bleeding, because ectopic pregnancy can be life-threatening.

\textbf{Contact your local emergency services for:}
\begin{itemize}
    \item Bleeding so heavy that you soak a pad or tampon every hour, or pass very large clots
    \item Dizziness, fainting, or collapse with bleeding
    \item Severe pain, especially one-sided lower abdominal pain, which may indicate an ectopic pregnancy
    \item Bleeding in pregnancy with pain, faintness, or shoulder-tip pain
    \item Heavy bleeding with fever, which may indicate infection
    \item Any bleeding that soaks through clothing and does not slow down
\end{itemize}

\section*{Diagnosis and Investigations}
The assessment of vaginal bleeding is guided by the woman's age, pregnancy status, and bleeding pattern.

\begin{itemize}
    \item \textbf{Pregnancy test:} always performed first in women of reproductive age
    \item \textbf{History:} bleeding pattern, cycle, contraceptives, medications, sexual history, and general health
    \item \textbf{Physical examination:} assessment of heart rate and blood pressure to check for significant blood loss, plus gentle pelvic examination to inspect the vulva, vagina, and cervix
    \item \textbf{Blood tests:} full blood count for anaemia, and blood grouping in those who are pregnant or bleeding heavily
    \item \textbf{Pelvic ultrasound:} the key test, examining the uterine lining, fibroids, polyps, and, in pregnancy, the location and viability of the pregnancy
    \item \textbf{Cervical screening and swabs:} for infection and cervical cancer screening where indicated
    \item \textbf{Endometrial biopsy:} recommended for women over 40 or those with postmenopausal bleeding to exclude endometrial cancer
    \item \textbf{Hysteroscopy:} direct visualisation of the uterine cavity when polyps, fibroids, or structural problems are suspected
\end{itemize}

\section*{Management}
Management depends entirely on the cause, the severity, and the woman's circumstances and fertility wishes. A calm explanation of the cause and the plan is important, because unexplained bleeding is frightening.

\begin{itemize}
    \item \textbf{Heavy menstrual bleeding:} tranexamic acid and NSAIDs during menstruation, the levonorgestrel intrauterine system, the combined pill, or other hormonal treatments; surgery such as endometrial ablation or hysterectomy is an option for severe cases
    \item \textbf{Fibroids and polyps:} hormonal treatment or surgical removal (myomectomy, polypectomy)
    \item \textbf{Hormonal imbalance:} cycle regulation with hormonal treatment and management of the underlying cause
    \item \textbf{Miscarriage:} expectant management, medicines to complete the miscarriage, or surgery to empty the uterus
    \item \textbf{Ectopic pregnancy:} urgent treatment with medicines or surgery to remove the ectopic pregnancy
    \item \textbf{Infection:} antibiotics and treatment of any sexually transmitted infection, with partner treatment
    \item \textbf{Postmenopausal bleeding:} endometrial sampling to rule out cancer, followed by treatment of the cause
    \item \textbf{Bleeding on anticoagulants:} review of the medication with the treating doctor, never stopping it without advice
    \item \textbf{Acute heavy bleeding:} hospital admission, fluids, blood transfusion if needed, and sometimes urgent surgery or hormonal measures to stop the bleeding
\end{itemize}

\section*{Complications}
\begin{itemize}
    \item Anaemia and iron deficiency from chronic or heavy blood loss, causing fatigue, breathlessness, and reduced exercise tolerance
    \item Severe acute blood loss requiring transfusion
    \item Missed or delayed diagnosis of ectopic pregnancy, which can rupture and be fatal
    \item Missed or delayed diagnosis of endometrial or cervical cancer
    \item Chronic pain and interference with work, study, and quality of life
    \item Psychological distress and worry about what the bleeding means
\end{itemize}

\section*{Prognosis and Outcomes}
The prognosis depends on the cause. Most abnormal bleeding is benign and responds well to hormonal treatment, simple measures, or minor surgery. Heavy periods can almost always be improved substantially, and modern treatments preserve fertility where desired. Ectopic pregnancy is treatable but is an emergency, and outcomes are excellent when it is caught early. Postmenopausal bleeding is benign in the great majority of cases, and even when endometrial cancer is found, it is usually detected early because it causes bleeding, so cure rates are high. Women who are pregnant and bleeding have a roughly 50\% chance of the pregnancy continuing normally, which is why reassurance and follow-up are important. The key to good outcomes is prompt assessment and an accurate diagnosis, so no woman should delay in seeking review for unexpected bleeding.

\section*{Prevention and Self-Care}
\begin{itemize}
    \item Track your periods and seek review for any change in pattern, heaviness, or timing
    \item Report any bleeding after sex, between periods, or after menopause promptly
    \item Use contraception effectively if you do not want to become pregnant, since early pregnancy is a common cause of unexpected bleeding
    \item Maintain a healthy weight and manage conditions such as thyroid disease and diabetes
    \item Take iron if advised for heavy periods, and eat iron-rich foods
    \item Attend cervical screening when it is due
    \item Seek medical review before starting or stopping any medication, including anticoagulants
    \item If you are pregnant, report any bleeding to your midwife or doctor, even if it is light
\end{itemize}

\section*{Sources and Further Reading}
The following organisations provide reliable information on vaginal bleeding and its causes.

\begin{itemize}
    \item Healthdirect Australia. \textit{Vaginal bleeding: causes and when to seek help.} https://www.healthdirect.gov.au/
    \item NHS. \textit{Vaginal bleeding and abnormal uterine bleeding.} https://www.nhs.uk/
    \item Mayo Clinic. \textit{Vaginal bleeding: when to see a doctor.} https://www.mayoclinic.org/
    \item MSD Manual. \textit{Abnormal uterine bleeding.} https://www.msdmanuals.com/
    \item MedlinePlus. \textit{Vaginal bleeding between periods.} https://medlineplus.gov/
    \item Jean Hailes for Women's Health. \textit{Bleeding problems and heavy periods.} https://www.jeanhailes.org.au/
\end{itemize}
